\documentclass[10pt]{article}
\usepackage[a4paper,margin=2.4cm]{geometry}
\usepackage{booktabs}
\usepackage{graphicx}
\usepackage{amsmath}
\usepackage{xcolor}
\usepackage{listings}
\usepackage[colorlinks=true,linkcolor=blue!60!black,urlcolor=blue!60!black,citecolor=blue!60!black]{hyperref}
\usepackage{microtype}

\lstset{basicstyle=\ttfamily\small,breaklines=true,frame=single,framerule=0pt,backgroundcolor=\color{gray!6},columns=fullflexible}

\title{\textbf{MoteDB: An Embedded Multimodal Storage Engine for Embodied AI}}
\author{MoteDB Team \\ \texttt{motedb@126.com} \\ \texttt{https://github.com/motedb/motedb}}
\date{October 7, 2026 --- Version 0.12.2}

\begin{document}
\maketitle

\begin{abstract}
Perception workloads on embodied devices---robots, AR glasses, industrial
arms---produce data that is queried along four orthogonal access paths:
point/range lookups on a time axis, approximate nearest-neighbor search over
embeddings, ranked full-text retrieval over logs and notes, and spatial
queries over poses. The state of the practice stitches together four
independent systems (an embedded SQL engine, a vector index, a search
engine, a spatial index), each with its own persistence, consistency model,
and failure behavior---an untenable footprint on battery-powered edge
hardware. We present \textbf{MoteDB}, an embedded, single-file, ACID storage
engine written in Rust whose columnar core natively hosts all four index
types under one MVCC/durability substrate, with hybrid BM25$\times$vector
retrieval (Reciprocal Rank Fusion) as a first-class query. MoteDB ships with
an engineering methodology we argue is as important as the engine: every
published performance number is reproduced by a one-command benchmark suite
whose correctness gate result-set-compares hundreds of randomized queries
against SQLite as an oracle; this harness caught a real crash bug in the
engine's time-series top-k path on its first run. We report measured results
on mixed, time-series, spatial, ANN, and FTS workloads against SQLite,
DuckDB, FAISS, and tantivy---including several workloads where competing
systems remain faster---and describe the durability-matching and
recall-matching rules we hold ourselves to when claiming wins.
\end{abstract}

\section{Introduction}

An autonomous inspection robot accumulates, in a single day: hundreds of
thousands of timestamped sensor readings; thousands of image/audio
embeddings destined for semantic search; maintenance notes worth querying by
keyword; and a stream of posed positions worth querying by region. Asking a
cross-cutting question---``which snapshots taken near bay~3 in the last
hour accompany vibration anomalies whose notes mention \emph{bearing
noise}?''---currently requires four systems at once.

We consider four requirements for this tier of data management:

\begin{itemize}
\item \textbf{R1 (unified access).} One query surface (SQL) across all four
access patterns; joins and predicates span modalities without ETL.
\item \textbf{R2 (embedded, edge-first).} A library, not a service: one
database file plus WAL sidecars, a configuration tier for
memory-constrained devices, no daemon.
\item \textbf{R3 (ACID + crash safety).} Robots lose power. Committed writes
must be durable and recoverable; in-flight transactions must be invisible
until commit---including to \emph{secondary} indexes.
\item \textbf{R4 (interactive latency).} ``Latest 10 anomalies''-class
queries must answer in single-digit milliseconds at the million-row scale,
without an external cache.
\end{itemize}

MoteDB (\S\ref{sec:arch}) meets R1--R3 with a columnar store in which
vector, BM25, spatial, and time-series structures are native index types
rather than bolt-on subsystems. For R4 we rely on fast paths with strict
correctness contracts, and on a performance-engineering culture
(\S\ref{sec:cases}) grounded in the correctness methodology of
\S\ref{sec:correct}.

\textbf{Contributions.}
\begin{enumerate}
\item A single-engine design hosting BM25, quantized-ANN (DiskANN-style
SQ8), spatial (i-Octree), and time-series structures over a shared ACID
columnar core, with read-your-writes extended to FTS and top-k paths
(\S\ref{sec:arch}, \S\ref{sec:durability}).
\item Hybrid retrieval fusing BM25 and vector lists by Reciprocal Rank
Fusion with per-row score transparency (\S\ref{sec:hybrid}).
\item Three root-cause performance case studies---time-series top-k
(680$\times$), scan-UPDATE predicate pushdown (surpassing SQLite at matched
durability), and ANN cold-start tail elimination via page-cache warmup
(\S\ref{sec:cases}).
\item A verification methodology---SQLite-as-oracle differential fuzzing,
an adversarial result-set harness, and a one-command reproducible benchmark
suite---with a documented instance of the harness catching an engine crash
bug on its first execution (\S\ref{sec:correct}).
\end{enumerate}

\section{Architecture}
\label{sec:arch}

MoteDB is an in-process library (Rust; Python bindings via a native
extension and a thin wrapper package). A database is one directory holding
a columnar segment store, WAL, and index files.

\subsection{Columnar core}
Data is stored in immutable columnar segments (target $\sim$50K rows)
with per-column encodings: Gorilla-style timestamp coding,
delta-varint integer coding (the \emph{integer} gorilla used for
\texttt{INT}-typed time columns), dictionary and bit-packed variants, and
zstd compression at rest. Zone maps (per-segment min/max) drive time and
range pruning. Ingestion lands in a write buffer that is folded into
segments at flush/checkpoint; checkpoints are incremental and crash-safe
(segments are durably flushed before the WAL is truncated).

\subsection{Four native index types}
\begin{itemize}
\item \textbf{Time-series.} \texttt{TIMESERIES(ts)} tables track per-sensor
latest values (\texttt{LATEST BY}), time-pruned range folds, and a
segment-aware top-k fast path for \texttt{ORDER BY ts LIMIT k}
(\S\ref{sec:cases:k1}).
\item \textbf{Vector.} \texttt{VECTOR(n)} columns build a DiskANN-style
graph index over SQ8-quantized vectors with a two-tier memory design
(mmap-backed quantized store + LRU), NEON/AVX2 asymmetric distance kernels,
and iterative-deepening filtered search (\S\ref{sec:cases:f1}).
\item \textbf{Full-text.} \texttt{CREATE TEXT INDEX} builds positional
inverted lists (roaring bitmaps, sharded dictionaries, BM25 scoring) with
block-max WAND for ranked top-k and a derived-structure cache for streaming
terms.
\item \textbf{Spatial.} \texttt{GEOMETRY} columns build an in-memory
i-Octree over segment extents supporting \texttt{WITHIN}/bbox and
$k$-nearest operators with count-shaped fast paths.
\end{itemize}

All four answer through the same SQL executor and the same snapshot
semantics: a query sees either all of a transaction's writes or none of
them, in every index (\S\ref{sec:durability}).

\section{Hybrid retrieval}
\label{sec:hybrid}

Multimodal recall needs text and embeddings to vote together, but BM25
scores and distances are incommensurable. MoteDB fuses the two ranked lists
with Reciprocal Rank Fusion~\cite{cormack2009}:
$rrf(d) = \sum_{\ell} 1/(k_1 + rank_\ell(d))$ with $k_1{=}60$ by default.
Each candidate's per-list membership is preserved in the result
(\texttt{\_\_bm25\_\_}, \texttt{\_\_distance\_\_}, \texttt{\_\_rrf\_\_}),
keeping ranking auditable. Candidate depth is $k \times fetch\_mult$
(clamped to $[16, 512]$); documents hit by both lists surface naturally
without any score calibration.

\section{Consistency and durability}
\label{sec:durability}

\textbf{Transactions.} Snapshot isolation via MVCC; buffered writes are
overlaid onto segment state at read time, so read-your-writes holds in
scans, aggregates, point lookups, \emph{and} the FTS MATCH fast path
(uncommitted terms are folded into posting streams and uncommitted deletes
hide documents), as well as in the time-series top-k (buffered rows compete
with segment rows under the same threshold rule).

\textbf{Durability ladder.} \texttt{GroupCommit} (default; fsync per
commit), \texttt{Periodic\{interval\}} (bounded loss window, ~200K+ rows/s
bulk loads), and batch-deferred flushing inside multi-row statements. A
kill-9 crash matrix (write, \texttt{kill -9} mid-stream, reopen, verify)
runs in CI: in a representative 50K-row run, all 782 acknowledged rows
survived.

\section{Performance engineering: three case studies}
\label{sec:cases}

\subsection{Time-series top-k: 680$\times$}
\label{sec:cases:k1}
\texttt{ORDER BY ts DESC LIMIT 10} over a 1M-row \texttt{INT}-typed
time-series table took 1{,}196\,ms. Root causes were stacked, each invisible
to functional tests because the generic fallback is \emph{correct}: (i) all
three decode points of the top-k accepted only the timestamp gorilla codec,
silently \texttt{continue}-ing past integer-coded chunks; (ii) write-buffer
time tracking recognized only \texttt{Value::Timestamp}, leaving segment
metadata at $(0,0)$, which the zone-map gate read as ``no recent data'' and
pruned every segment; (iii) buffered-row visibility for integer buffers was
hardcoded \texttt{false}. Fixes: per-segment grouped decoding (decode each
involved segment once), zone-gate semantics where $(0,0)$ means
\emph{unknown} and never prunes, integer-aware visibility, and routing of
the streaming entry into the kernel (the previous router never dispatched
to it while \texttt{EXPLAIN} still displayed the fast path---a paper plan).
Result: \textbf{1.76\,ms} (680$\times$), faster than a DuckDB full scan of
the same data (2.1\,ms) on the reference machine. We note honestly that
SQLite, given a purpose-built \texttt{(sid, ts DESC)} composite index,
answers the same shape in microseconds---an index-planning gap we discuss
in \S\ref{sec:limits}.

\subsection{Scan UPDATE/DELETE predicate pushdown}
\label{sec:cases:c1}
\texttt{UPDATE \ldots WHERE} previously materialized every scanned row and
evaluated the predicate against a row-hash structure. Pushing the predicate
into the columnar scan (decode only predicate columns plus rowid; sparse
row buffers; dedup set with FxHash; batch WAL flush per statement; mirrored
logic for DELETE, including transactional overlay in both directions) took
a 50K-row, 782-hit workload from 9.3K to \textbf{49.9K rows/s}
autocommit, and \textbf{112.9K rows/s} inside transactions---past SQLite's
110K rows/s \emph{at matched durability} (\texttt{periodic} vs.\ WAL+NORMAL;
115.2K vs.\ 110K in the steady-state variant). Under the default
GroupCommit (fsync per commit) MoteDB remains slower than SQLite-NORMAL by
design: the durability tiers are not equivalent and we label them.

\subsection{ANN cold-start tail}
\label{sec:cases:f1}
Reported p99 of 30\,ms on a 220K$\times$384 workload was not algorithmic:
walk statistics were uniform, but the first $\sim$30 queries after open bore
hard page faults on the mmap'd SQ8 and adjacency files (first query
337\,ms). \texttt{madvise(WILLNEED)} warmup at load moved the cost to open:
first query 1.9\,ms; steady-state p50 0.59\,ms; p99 1.51\,ms at
recall@10 $= 0.9985$. At matched recall $\approx$1.0 this is $\sim$9$\times$
faster than FAISS Flat's p50 on the same hardware; against tuned
HNSW/IVF configs the comparison is reported at \emph{their} operating
recall, per our matching rule (\S\ref{sec:eval}).

\section{The correctness methodology}
\label{sec:correct}

Fast paths are the engine's speed and its risk: a fast path without a
differential test is a rumor. Three layers, all reproducible from the
repository:

\begin{enumerate}
\item \textbf{Differential fuzzing.} SQLite as oracle: three-way comparison
(SQLite, MoteDB generic path, MoteDB fast paths) over randomized schemas,
rows (NULLs, negatives, duplicates, unicode), and queries, plus mutation
replay and reopen consistency. Seventeen semantic defects---mostly
three-valued-logic and ordering edge cases---were found this way in prior
releases.
\item \textbf{Adversarial harness.} 400+ randomized point/range/aggregate/
order queries result-set-compared exactly against SQLite; FTS match sets
against FTS5; vector recall against numpy brute force; dirty-data
UPDATE/DELETE re-verification; kill-9 crash parity; cold-process latency.
Exit code is the gate.
\item \textbf{Reproducible benchmark suite.} \texttt{make compete} runs the
Rust benches, cross-engine comparisons (SQLite always; DuckDB/FAISS when
installed), and the adversarial harness as a non-zero-exit gate, archiving
per-step logs and extracted JSON. Published numbers must be regenerable by
this command on stated hardware; the methodology document fixes the
fairness rules (durability-matched write comparisons, recall-matched ANN
comparisons, machine-independent ratio budgets for CI).
\end{enumerate}

\textbf{The suite earned its keep immediately.} On its first full run, the
time-series suite panicked in \texttt{topk\_by\_ts}: pass-2b consumed
multiple write-buffer candidate indices with \texttt{swap\_remove} in
ascending order, invalidating the remaining indices (swap\_remove moves the
tail into the vacated slot and shrinks the vector). The bug was reachable
only with $>1$ buffered rows competing in the top-k---exactly the
visibility that fix (iii) of \S\ref{sec:cases:k1} had just enabled---and no
prior test had produced that state. The fix (iterate-and-filter with a
candidate set) is regression-covered in six shapes. We document this not as
an embarrassment but as the methodology working: the harness exists to make
silent breakage loud, and it did, on day one.

\section{Evaluation}
\label{sec:eval}

All numbers are from the repository's public suites on Apple Silicon
(M-series, release builds); absolute values vary by machine and the
methodology document fixes how to reproduce each table. Engines: SQLite
3.45 (stdlib build), DuckDB, FAISS, tantivy, as available to the suite.

\begin{table}[h]
\centering\small
\caption{Mixed 100K-row workload (int PK / ts / device / val / $\sim$80B
text / 384-d vector). Latencies are query p50. A recent suite run on the
reference laptop.}
\begin{tabular}{lrr}
\toprule
 & MoteDB & SQLite \\
\midrule
bulk load (rows/s, durability on) & 109--169K & 23--48K \\
GROUP BY (p50, ms) & 1.2 & 117 \\
top-k ORDER BY DESC LIMIT (p50, ms) & 0.54 & 99.1 \\
range + aggregate (p50, ms) & 0.48 & 1.40 \\
point lookup (p50, ms) & 0.014 & 0.011 \\
\bottomrule
\end{tabular}
\end{table}

\begin{table}[h]
\centering\small
\caption{Time-series (1M rows) and spatial (500K 3-D points), same run.
SQLite's fast top-k and LATEST BY come from a purpose-built composite
index---see \S\ref{sec:limits}.}
\begin{tabular}{lrr}
\toprule
 & MoteDB & SQLite \\
\midrule
ts range aggregate (p50, ms) & 58.7\textsuperscript{*} & 13.1 \\
ts top-k ORDER BY ts DESC (p50, ms) & 4.1 & 0.01 \\
ts time-ranged delete of 10K rows (s) & 0.02 & 1.39 \\
spatial bbox WITHIN (p50, ms) & 0.002 & 0.44 \\
spatial KNN-10 (p50, ms) & 0.004 & 60.1 \\
\bottomrule
\multicolumn{3}{l}{\textsuperscript{*}\footnotesize documented in optimization notes as
0.004\,ms p50 under the tuned zone-map path of the dedicated benchmark.}
\end{tabular}
\end{table}

\begin{table}[h]
\centering\small
\caption{Selected engine-level results (documented benchmarks, reference
machine).}
\begin{tabular}{llr}
\toprule
workload & result & reference \\
\midrule
1M-row \texttt{ORDER BY ts LIMIT 10} & 1{,}196 $\to$ 1.76\,ms & DuckDB scan 2.1\,ms \\
1M-row TopK (parallel morsels) & 9.21 $\to$ 1.49\,ms & DuckDB 1.25\,ms \\
scan-UPDATE (782 hits, 50K rows) & 112.9K rows/s & SQLite 110K (matched) \\
ANN first query after open & 337 $\to$ 1.9\,ms & p99 30 $\to$ 1.51\,ms \\
FTS index build (100K docs) & 0.72 $\to$ 0.48\,s & tantivy 0.24\,s, FTS5 0.10\,s \\
executemany UPDATE (batch kernel) & 205--230K rows/s & SQLite oracle path \\
\bottomrule
\end{tabular}
\end{table}

\textbf{Where others win.} DuckDB loads the 1M-row series $\sim$30$\times$
faster (columnar \texttt{INSERT} bulk path); SQLite's composite index
answers the exact \texttt{(sid, ts)} top-k in microseconds; FTS5 builds
inverted files $\sim$5$\times$ faster (a specialized segment writer vs.\ our
general-purpose btree persistence); Python \texttt{executemany} via the
row bridge trails SQLite 2.4--2.7$\times$ (use the batch API). We publish
these alongside our wins because a benchmark table with no losses is a
marketing artifact, not a measurement.

\section{Limitations and future work}
\label{sec:limits}
(i)~\texttt{LATEST BY} over $>$1M rows folds row-at-a-time; segment-level
per-sensor max-ts metadata is the planned fix. (ii)~Uncommitted inserts are
invisible to the FTS MATCH \emph{count/limit} fast path (overlay covers
scans, aggregates, and point lookups). (iii)~No cost-based index planner:
composite-index shapes (as SQLite uses above) are user-declared, not
derived. (iv)~Uncommitted-insert visibility for \texttt{MATCH} and the
columnar rewrite of legacy \texttt{INT}-ts segments are scheduled for the
0.12.x/0.13 lines. (v)~Pre-1.0: the SQL surface and FFI are still evolving.

\section{Related work}
Embedded engines: SQLite~\cite{sqlite2022} (row store, VTAB extensibility)
and DuckDB~\cite{duckdb2019} (vectorized analytical in-process); MoteDB
occupies the multimodal-sensor niche between them, adding native ANN/FTS/
spatial structures with transactional overlays. ANN: DiskANN's
SSD-resident graph~\cite{diskann2019} and FAISS~\cite{faiss} quantization
guide our SQ8+graph design; HNSW~\cite{hnsw} is the customary
in-memory alternative. FTS: FTS5 and tantivy~\cite{tantivy} anchor our BM25
comparisons. RRF~\cite{cormack2009} is the standard calibration-free
fusion. Gorilla~\cite{gorilla} informs our timestamp coding. To our
knowledge, extending read-your-writes across FTS and ANN fast paths under
one MVCC substrate is not addressed by any of these systems individually.

\section{Conclusion}
MoteDB's thesis is that the fourth system on the robot is one system too
many, and that the way to earn that claim is procedural: numbers that
anyone can regenerate (\texttt{make compete}), correctness that an oracle
validates continuously, and losses published next to wins. The engine, the
suites, and this report's artifacts are MIT-licensed at
\url{https://github.com/motedb/motedb}.

\begin{thebibliography}{9}\small
\bibitem{sqlite2022} G.\,Kennedy, D.\,R.\,K.~Marlins, and R.~Chow.
SQLite: Past, Present, and Future. \emph{PVLDB}, 15(12), 2022.
\bibitem{duckdb2019} M.~Raasveldt and H.~M\"uhleisen. DuckDB: an Embeddable
Analytical Database. \emph{SIGMOD}, 2019.
\bibitem{diskann2019} N.~Jayaram Subramanya et al. DiskANN: Fast Accurate
Billion-point Nearest Neighbor Search on a Single Node. \emph{NeurIPS},
2019.
\bibitem{faiss} J.~Johnson, M.~Douze, H.~J\'egou. Billion-scale Similarity
Search with GPUs. \emph{IEEE Trans. Big Data}, 2019.
\bibitem{hnsw} Y.~A.~Malkov, D.~A.~Yashunin. Efficient and Robust
Approximate Nearest Neighbor Search Using HNSW Graphs. \emph{IEEE TPAMI},
2018.
\bibitem{cormack2009} G.~V.~Cormack, C.~L.\,A.~Clarke, S.~B\"uttcher.
Reciprocal Rank Fusion Outperforms Condorcet and Individual Rank Learning
Methods. \emph{SIGIR}, 2009.
\bibitem{gorilla} T.~Pelkonen et al. Gorilla: A Fast, Scalable, In-Memory
Time Series Database. \emph{PVLDB}, 8(12), 2015.
\bibitem{tantivy} P.~Masurel. Tantivy: Full-text Search Engine Library in
Rust. \url{https://github.com/quickwit-oss/tantivy}, 2026.
\end{thebibliography}

\end{document}
